\section{Introduction}
RAG systems augment a language model with records selected from an external datastore \cite{lewis2020rag}. In a healthcare setting, this creates a security question that is more fundamental than whether a model eventually prints protected information: \emph{did the system retrieve protected information into model context at all?} If authorization is enforced only through instructions inside the prompt, the model may be told not to reveal a record after the record has already crossed the security boundary.

Recent work has begun integrating permissions and metadata filtering into RAG pipelines. Chen \emph{et al.} demonstrated access control over patient profiles in a 40-profile proof of concept \cite{chen2025sac}; Jeong and Lee proposed IAM-backed permission-aware filtering for multi-resource environments \cite{jeong2025permission}; and Chen and Taipalus combined metadata filtering with fine-grained role-based access control \cite{chen2026metadata}. In parallel, extraction attacks show that RAG datastores can be deliberately induced to leak through instruction following and adaptive querying \cite{qi2025spill,dimaio2025pirates}. These lines of work motivate a careful distinction between retrieval-time exposure and generation-time leakage, and between a structural access-control boundary and heuristic attack detection.

This paper evaluates three architectures under the same synthetic healthcare benchmark: (i) a prompt-only baseline that retrieves globally and places authorization policy in the system prompt, (ii) deterministic pre-retrieval ACL filtering, and (iii) the same ACL boundary with an additional risk-triage controller. Our central claim is deliberately narrow: \emph{authorization should constrain the retrieval candidate set before sensitive content is assembled into model context}. We do not claim clinical efficacy, model-level privacy, or authentication correctness; principal resolution is assumed to have occurred before the evaluated pipeline begins.

The study makes four contributions. First, it separates Unauthorized Context Exposure Rate (\ucer) from a narrow canary Unauthorized Disclosure Rate (\udr), preventing a low output-leak rate from masking a complete retrieval-boundary failure. Second, it resolves an initially observed utility gap with a paired, controlled $k=1$ experiment that equalizes retrieval composition. Third, it compares fuzzy risk triage with a transparent deterministic-rule baseline using label-independent telemetry and lexically disjoint held-out suspicious prompts; the negative generalization result constrains what can credibly be attributed to risk scoring. Fourth, it adds a secondary, nested 1K--100K systems study that separates global-corpus growth from authorization-scope growth while fitting a single retrieval representation on the 100K master corpus.
